If $ax=0$ and $by=0$ with $a,b\ne0$, then $ab(x+y)=0$ and $a(rx)=0$ for every $r\in A$. Since $ab\ne0$, $T(M)$ is a submodule.

\textbf{(i)} If $a(x+T(M))=0$ with $a\ne0$, then $ax\in T(M)$, so $bax=0$ for some $b\ne0$. Hence $x\in T(M)$.

\textbf{(ii)} If $ax=0$, then $af(x)=f(ax)=0$.

\textbf{(iii)} Write the maps as $f$ and $g$. Injectivity on $T(M')$ follows from that of $f$. If $x\in T(M)$ and $g(x)=0$, write $x=f(y)$ and choose $a\ne0$ with $ax=0$. Then $f(ay)=0$, so $ay=0$ and $y\in T(M')$. The reverse inclusion in the kernel is immediate.

\textbf{(iv)} If $ax=0$ with $a\ne0$, then $1\otimes x=(1/a)\otimes ax=0$. Conversely, $K$ is the directed union of the submodules $A(1/s)$ for $s\ne0$. Since tensor products commute with directed limits, $1\otimes x=0$ in $K\otimes_A M$ implies its vanishing in $A(1/s)\otimes_A M$ for some $s\ne0$. Under the isomorphism $(a/s)\otimes m\mapsto am$, this element maps to $sx$. Thus $sx=0$.
